\chapter{Every Other Kitchen in Town}

\epi{A restaurant guide to the dark sector.\\ Stars awarded for flavour. Nobody gets the fourth star.}

\section{The crawl}

Once the main recipe was on record, the obvious objection was: you only tried \emph{your} version of \LCDM. What about self-interacting dark matter? Fuzzy dark matter? The simulations themselves? Superfluids, dipoles, emergent gravity? So, over a single long day, the paper went on a restaurant crawl. Amendments A16 through A48 each took one more kitchen and graded it on the same frozen gates: the distance to the law (rms $\le0.02$, 95th-percentile galaxy $\le0.04$, worst bin $\le0.02$, unseen galaxies $\le0.025$), and whether the scale ties to $\rhoDE$, including, from A29 on, a tracking gate where dark energy evolves.

Here's the guide. Ratings are mine and carry no physics; the numbers are the paper's.

\section{More knobs on the same stove}

\textbf{A16, the knob blocks.} Four new blocks of knobs added to the fourteen-knob recipe, keeping whatever helps. The best set has 15 knobs and reaches $0.0263\dex$; with widened bounds $0.0252$, at a median halo mass $0.64\dex$ (4.4 times) above the stellar-to-halo relation. On unseen galaxies it gives $0.0267$, and the per-galaxy 95th percentile, $0.053$, blocks it. This is the closest single recipe in the whole paper.

\textbf{A25, the kitchen sink.} 28 \LCDM\ ingredients stacked into one global recipe: the 14 feedback knobs, the A16 blocks, two published halo-response forms, a self-interacting core, adiabatic contraction, separate settings for gas-rich and bulge galaxies, the slope and scatter of the stellar-to-halo relation. Plot distance against knob count and the curve \emph{flattens}, near $0.022\dex$; 28 knobs give $0.0222$. Widen the bounds and it reaches $0.0204$, but $0.0370$ on unseen galaxies, at a plausibility cost of 1.29 range-widths outside the literature. Its $\kappa_{\rm eff}$ of 0.497 is inherited from the target, and its scale still rises about three times by $z=2.33$.

\begin{knobbox}{The plateau}
Adding knobs stops helping at about $0.022\dex$. That's the most interesting single number in this chapter. A model that could reach the law if only it had enough freedom would keep descending as knobs are added. A plateau means the remaining gap is a shape the ingredients can't make, not a level they can't reach.
\end{knobbox}

\textbf{A26, the same recipe in other universes.} The frozen 14-knob recipe, put into universes with different $\rhoDE$, moves with slope $+0.187$ to $+1.453$ depending on what's held fixed, against 0.5. By $z=2.33$ its scale is 3.02 to 4.55 times today's. You'll see that bar again in the figure of Chapter~13.

\section{Other kinds of dark matter}

\begin{description}[leftmargin=1em,style=nextline]
\item[Wave (fuzzy) dark matter, A17.] A solitonic core on an NFW halo, from an ultralight boson. The boson mass alone gives $0.149\dex$; six knobs give $0.092$. The fits improve by \emph{switching the soliton off}: at six knobs it survives in none of the 171 galaxies, and the soliton's acceleration scale grows with halo mass.
\item[Hydro-calibrated halo responses, A18.] Published formulas fitted to simulations (Di Cintio and collaborators 2014, coreNFW, and the FIRE-2 form of Lazar and collaborators 2020) sit $0.116$ to $0.177\dex$ from the law; with three knobs each, $0.073$ to $0.084$, wanting lighter halos.
\item[Self-interacting dark matter, A19 and A41.] Velocity-dependent cross-section: $0.122\dex$, $0.080$ with three halo knobs. Per-galaxy halo masses and concentrations at $\sigma/m=0.5$\,cm$^2$\,g$^{-1}$, in the style of Ren and collaborators (2019): $0.0487$ with 343 parameters, 77 at a bound, and $0.110$ on unseen galaxies. On its collapsed branch, $0.179$; with each galaxy free to choose its branch, $0.047$ with 514 parameters but $0.110$ unseen.
\item[Warm and mixed dark matter, A19 and A41.] $0.116$ to $0.120$.
\item[Condensate dark matter, A41.] A fixed core radius gives $0.129$, with no core solution at all in 83 of 171 galaxies.
\item[Decaying dark matter, primordial black holes, dark radiation, a Yukawa force, A42.] Best $0.0717$, with every fitted particle parameter beyond published limits. The black holes, for example, come out at $4\times10^5\Msun$, which would also overheat the discs.
\item[Dark mass that traces the gas, A43.] $0.41$ as published, $0.21$ with a stellar term, because bright discs hold too little gas.
\item[Halo flattening, tides, substructure, environment, A44.] Moves the best 28-knob recipe from $0.0222$ to $0.0221$. The plateau doesn't budge.
\end{description}

\section{The simulations themselves}

The fairest test of all would be real hydrodynamic simulations, which the paper can't run. Amendment A21 used their published results instead, as read from the papers. No measured simulation scale reaches the law. The closest is Magneticum, $0.0290\dex$, then MassiveBlack-II's faint galaxies at $0.0306$ and ZOMG at $0.0359$; EAGLE's raw fits sit $0.17$ to $0.19\dex$ away. Only round, order-of-magnitude figures of ``$10^{-10}\ms$'' pass. The EAGLE runs with varied $\Lambda$ report no acceleration relation, so the slope with $\rhoDE$ can't be tested from any simulation on record. Where redshift numbers exist, the scale rises into the past, by 1.305 to 1.428 times in ZOMG.

\begin{honestbox}
These simulation values come from the papers' summaries and are provisional. Simulations are the one kitchen this book can't fully inspect, and a future simulation could do better than anything graded here. If one does, and passes the five tests of Chapter~7, that would be a real result.
\end{honestbox}

\section{The near-MOND kitchens}

Several theories reproduce modified-dynamics phenomenology from something else. They're interesting because they're closest in spirit to my law.

\textbf{Superfluid dark matter, A22} (Berezhiani and Khoury, 2015). Dark matter condenses in galaxies into a superfluid whose phonons mediate an extra force on baryons. Its deep force, $\gbar+\sqrt{a\gbar}$, sits $0.0392\dex$ from the law. Setting its energy scale by $\Lambda^4=\rhoDE$ gives $\kappa_{\rm eff}=5.01$, ten times too big; reaching 0.5 needs its coupling set by hand, so $\kappa$ just moves into the coupling.

\textbf{Dipolar dark matter, A23} (Blanchet and Le Tiec, 2008). Its published cubic function gives the same $\gbar+\sqrt{a\gbar}$ and the same $0.0392$. Its free function \emph{can} take the law's shape exactly: a fourth-order expansion with three knobs reaches $0.0036$. But that's a fit of shape inside a reformulation of MOND, not a \LCDM\ mechanism, and its link to $\Lambda$ is stated as an order of magnitude and imposed.

\textbf{Emergent gravity, A24} (Verlinde, 2017). Best $0.0392$ with one knob. With today's $H_0$ it gives $\kappa_{\rm eff}=0.583$ and rises 3.51 times by $z=2.33$. With the de Sitter rate $\sqrt{\Lambda/3}$ it gives $\kappa_{\rm eff}=0.482$, slope 0.5 and no drift, but its effective scale spreads $0.115\dex$ between galaxies, and the choice of rate is imposed. It's the $\Lambda$ coincidence of Chapter~6 with $1/6$ in place of $1/2\pi$.

\begin{mathbox}{Three theories, one wrong shape}
Superfluid, dipolar and emergent gravity all reduce, in their plain forms, to
\[
g_{\rm sum}=\gbar+\sqrt{a\,\gbar},
\]
while the law is
\[
g_{\rm law}=\frac{\gbar}{2}+\sqrt{\frac{\gbar^2}{4}+a\,\gbar}.
\]
Both go to $\sqrt{a\gbar}$ in the deep regime ($\gbar\ll a$), and both go to $\gbar$ in the Newtonian regime, but at different rates. The excess over Newton is $g-\gbar\to\sqrt{a\gbar}$ for the sum, which keeps growing, and $g-\gbar\to a$ for the law, which saturates. At the transition, $\gbar=a$:
\[
g_{\rm sum}=2a,\qquad g_{\rm law}=\tfrac12a+\sqrt{\tfrac54}\,a=1.618\,a,
\]
the golden ratio, as it happens. The gap there is $\log_{10}(2/1.618)=0.092\dex$. Averaged over SPARC's points, it's the $0.0392\dex$ that all three theories share. Same functional form, same miss: the transition has the wrong shape.
\end{mathbox}

\fig{sum_vs_law}{Three theories, one shape. Left: the boost $g/g_{\rm bar}$ for the law and for the sum form $g_{\rm bar}+\sqrt{ag_{\rm bar}}$. Right: their difference, which peaks near the transition, where the law gives the golden ratio and the sum form gives 2.}


\section{Fitting the observed curves directly}

Amendment A28 changed the question. Forget imitating the law; fit each galaxy's \emph{observed} rotation curve with mass-to-light ratio, distance and inclination free within published priors, and compare cored \LCDM\ halos with my law (with its dark energy tie).

\textbf{Here the rival wins.} Cored halos fit better than the tied law even after paying the BIC rent: total BIC 17804.6 against 19744.2, and they beat the law in 154 of 171 galaxies.

But the fitted galaxies, put back through the radial acceleration relation, give a relation $+0.056$ to $+0.103\dex$ off the tied $a_0$, with $0.125$ to $0.137\dex$ of scatter. So they reproduce the data \emph{without} reproducing the law. A better fit is not a derivation, and it isn't proof for either side. The paper's verdict: \textbf{the swoop that works is fitting the curves galaxy by galaxy; the one that fails is deriving the law and its tie.}

\section{Evolving dark energy and the rival literature}

Amendments A29 to A40 added the sharpest gate: where dark energy evolves, the scale must follow $\sqrt{\rhoDE(z)}$ to within $0.05\dex$ at $z=1$, 1.5 and 2.33.

\textbf{\LCDM\ halos follow the wrong thing.} In evolving backgrounds, halo scales follow the expansion rate and the matter density, not dark energy. Their slope with $\rhoDE$ runs from $-0.124$ to $+0.026$, and none of 315 $w_0$--$w_a$ backgrounds passes with the 14-knob recipe frozen (A29). Tuning feedback and dark energy together reaches tracking only through fitted powers of $H(z)$ with no named process, $H^{+3.0}$ on the stellar-to-halo relation and $H^{+4.2}$ on feedback (A32), which then distort the curves at high redshift and fail on unseen galaxies.

\textbf{The best rival line in the literature (A34).} The cooling scale of Kaplinghat and Turner (2002), one of the earliest \LCDM\ explanations of a MOND-like scale, put inside the law's formula and multiplied by one factor, 0.398, set by hand, matches today's curves exactly: $0.0152\dex$, $0.0156$ unseen. The shape is borrowed from the law, and the scale is tuned. It follows the matter density, not $\rhoDE$, and rises 1.42 times by $z=2.33$. Disc-formation and halo-response papers miss by 0.12 to $0.22\dex$.

\textbf{Screened fields (A30).} A three-knob k-mouflage fit with no halo reaches the law's shape, $0.0086\dex$, because its switch lies below almost every point and it acts as a power law. It adds $61.6\,a_0$ at Saturn, which the Cassini spacecraft would have noticed, and its scale doesn't depend on $\rhoDE$. Chameleon, symmetron and coupled quintessence fail.

\textbf{The logotropic fluid (A36)} of Chavanis carries the slope 0.5 natively, the only route that does without help, but its halo pull caps at $0.298\,a_0$ where the law needs $0.618\,a_0$, and its scale stays flat as dark energy evolves.

\textbf{De Sitter and quantised inertia (A37)} share one shape, $0.0189\dex$ with a free constant, but fail the worst bin at $0.0273$, with $\kappa_{\rm eff}=0.651$. Dark matter tied to baryons (A38): $0.072$ at best. Clustering dark energy (A39): at most $0.1\,a_0$ in halos. Correlated halo scatter (A40): best $0.154$.

\section{Modified gravity and systematics}

The last four amendments asked whether other modified-gravity theories, or plain measurement error, could produce the law with its tie.

\textbf{Modified gravity (A45, A46).} Vainshtein-screened DGP and Galileon gravity leave galaxies Newtonian, $0.496\dex$ from the law (or $0.142$ with a \LCDM\ halo). Beyond-Horndeski theories have no acceleration scale. In relativistic MOND theories (TeVeS, AeST, bimetric MOND) $a_0$ is a constant written into the Lagrangian, and where the interpolation is free it can be chosen to be the law, which is a choice, not a derivation. TeVeS's standard function is the sum form, $0.10$ at its usual $a_0$. MOG misses by $0.377$, conformal gravity by $0.136$.

\textbf{Measurement systematics (A47, A48).} Is the law just what ordinary halos look like through imperfect measurements? Non-circular motions, bars, pressure support, disc thickness, beam smearing, selection, mass-to-light trends and gas uncertainties, applied to \LCDM\ mocks at literature amplitudes, reach at best $0.122\dex$. Every mock keeps an excess of $+0.13$ to $+0.16\dex$ in the lowest-acceleration bin. A selection weight towards weak halos reaches the scale, $\kappa_{\rm eff}=0.518$, but not the shape.

\begin{honestbox}[A systematic that cuts my way, and against me]
The same corrections, applied to the \emph{real} SPARC data, raise the law's own fitted $a_0$ by up to $+0.101\dex$. Allowed mass-to-light and gas changes move it by $-0.174$ to $+0.157\dex$. In other words, my one knob, $\kappa=\tfrac12$, carries a systematic uncertainty of order $\pm0.15\dex$, roughly 30 to 40 per cent either way. That's a big reason why $\kappa$ is described as fitted and why the data can't single out $\tfrac12$ (Chapter~10).
\end{honestbox}

\section{The guide, summarised}

\fig{kitchens_map}{Every kitchen on one map: best distance to the law against number of knobs. \LCDM\ recipes (dark) flatten out near 0.022\,dex. The kitchens that get below the target (green) do it by putting the law's shape in by choice. None ties the scale to dark energy.}


\begin{table}[h]\centering\small
\begin{tabular}{@{}lrrl@{}}
\toprule
kitchen & knobs & best dex & tie to $\rhoDE$?\\
\midrule
\LCDM\ feedback, A16 & 15 & 0.0252 & no\\
kitchen sink, A25 & 28 & 0.0222 & inherited; rises $3\times$\\
Kaplinghat--Turner, A34 & 1 + law shape & 0.0152 & no; rises $1.42\times$\\
simulations, A21 & (published) & 0.0290 & untestable\\
superfluid, A22 & 1 & 0.0392 & $\kappa$ 5.0 natively\\
dipolar, A23 & 3 & 0.0036 & imposed\\
emergent, A24 & 1 & 0.0392 & imposed rate\\
k-mouflage, A30 & 3 & 0.0086 & no; fails Saturn\\
logotropic, A36 & native & caps at $0.3\,a_0$ & slope yes; flat in $z$\\
SIDM per galaxy, A19 & 343 & 0.0487 & no\\
coreNFW on data, A28 & per galaxy & beats law on BIC & law not reproduced\\
\bottomrule
\end{tabular}
\caption{Every other kitchen. None gets the law \emph{and} the tie natively.}
\end{table}

Nobody gets the fourth star. Some kitchens reach the law's shape (dipolar, k-mouflage) by choosing a free function. Some reach today's curves exactly by borrowing the law's form and tuning a scale (Kaplinghat--Turner). One beats the law outright on observed curves (per-galaxy cored halos) without reproducing the law. And not one of them, on its own terms, makes the acceleration scale follow the density of dark energy.

Which brings us to a theory with rather more freedom than any kitchen in town.
